% Part: incompleteness
% Chapter: incompleteness-provability
% Section: first-incompleteness-thm

\documentclass[../../../include/open-logic-section]{subfiles}

\begin{document}

\olfileid{inc}{inp}{1in}

\olsection{మొదటి అసంపూర్ణతా సిద్ధాంతం}

ఇప్పుడు మొదటి అసంపూర్ణతా సిద్ధాంతానికి గ్యోడెల్ ఇచ్చిన అసలు నిరూపణను
వివరించవచ్చు. అంకగణిత భాషను విస్తరించే ఒక భాషలో గణనీయంగా
స్వీకృతీకరించబడిన ఏ సిద్ధాంతమైనా $\Th{T}$గా తీసుకోండి; ఆ $\Th{T}$ అనేది
$\Th{Q}$ స్వీకృతాలను కలిగి ఉండనివ్వండి.
దీని వల్ల, ప్రత్యేకించి, $\Th{T}$ గణనీయ ప్రమేయాలకు, సంబంధాలకు ప్రాతినిధ్యం
వహిస్తుంది.

సూత్రాలనూ నిరూపణలనూ సంఖ్యలుగా సమంజసంగా సంకేతీకరిస్తే,
$\Prf[\Th{T}](x,y)$ సంబంధం గణనీయమైనదని మనం వాదించాం; ఇక్కడ
$\Prf[\Th{T}](x,y)$ అంటే $x$ అనేది $\Th{T}$లో గ్యోడెల్ సంఖ్య~$y$ గల
\tetoken{సూత్రపు}{!!{formula}} \tetoken{ఒక వ్యుత్పత్తికి}{!!a{derivation}} గ్యోడెల్ సంఖ్య అని అర్థం. నిజానికి,
గ్యోడెల్ దృష్టిలో పెట్టుకున్న ప్రత్యేక సిద్ధాంతానికి, తన పత్రంలోని 45
ప్రమేయాలు, సంబంధాల జాబితాను వాడి ఈ సంబంధం ఆదిమ పునరావర్తనమని ఆయన
చూపగలిగాడు. ఆయన ఎంచుకున్న ప్రత్యేక~$\Th{T}$కు, 45వ సంబంధం $x B y$ అనేది
$\Prf[\Th{T}](x,y)$యే. గ్యోడెల్ తన పత్రంలో ``పునరావర్తన'' అనే పదం
వాడిన చోట, ఇప్పుడు మనం ``ఆదిమ పునరావర్తన'' అని అంటామని గుర్తుంచుకోండి.

$\Prf[\Th{T}](x,y)$ గణనీయమైనది కాబట్టి, అది $\Th{T}$లో ప్రాతినిధ్యం
చేయదగినది. దానికి ప్రాతినిధ్యం వహించే సూత్రాన్ని $\OPrf[\Th{T}](x,y)$తో
సూచిస్తాం. $\OProv[\Th{T}](y)$ అనేది
$\lexists[x][\OPrf[\Th{T}](x,y)]$ అనే సూత్రంగా ఉండనివ్వండి. ఇది గ్యోడెల్
జాబితాలోని 46వ సంబంధం $\fn{Bew}(y)$ను వర్ణిస్తుంది. గ్యోడెల్ చెప్పినట్లు,
``పునరావర్తనమని ప్రకటించలేని'' ఏకైక సంబంధం ఇదే. బహుశా ఆయన ఉద్దేశం ఇది:
నిర్వచనం నుంచి అది గణనీయమని స్పష్టం కాదు; తదుపరి పరిణామాలు నిజానికి అది
గణనీయం కాదని చూపిస్తాయి.

$\Th{T}$ అనేది $\Th{Q}$ను కలిగి ఉన్న \tetoken{స్వీకృతీకరించదగిన}{!!{axiomatizable}}
సిద్ధాంతం అనుకుందాం. అప్పుడు $\Prf[\Th{T}](x, y)$ నిర్ణయించదగినది; కాబట్టి
\tetoken{ఒక సూత్రం}{!!a{formula}}~$\OPrf[\Th{T}](x, y)$ ద్వారా $\Th{Q}$లో ప్రాతినిధ్యం
చేయదగినది. పైన వివరించిన సూత్రమే $\OProv[\Th{T}](y)$గా ఉండనివ్వండి.
స్థిరబిందు ఉపసిద్ధాంతం ప్రకారం, $\Th{Q}$ (అందువల్ల $\Th{T}$ కూడా) కింది
దాన్ని \tetoken{నిగమించే}{!!{derive}s} సూత్రం $!G_\Th{T}$ ఉంటుంది:
\begin{equation}
\ollabel{eqn:qpf}
!G_\Th{T} \liff \lnot \OProv[\Th{T}](\gn{!G_\Th{T}}).
\end{equation}
$!G_\Th{T}$ స్థూలంగా ``$!G_\Th{T}$ అనేది $\Th{T}$లో
\tetoken{వ్యుత్పాదించదగినది}{!!{derivable}} కాదు'' అని చెబుతుందని గమనించండి.

\begin{lem}\ollabel{lem:cons-G-unprov}
$\Th{T}$ అనేది $\Th{Q}$ను విస్తరించే అవైరుధ్యమైన,
\tetoken{స్వీకృతీకరించదగిన}{!!{axiomatizable}} సిద్ధాంతమైతే, $\Th{T} \Proves/ !G_\Th{T}$.
\end{lem}

\begin{proof}
$\Th{T}$ అనేది $!G_\Th{T}$ను \tetoken{నిగమిస్తుందని}{!!{derive}s} అనుకుందాం. అప్పుడు
\tetoken{ఒక వ్యుత్పత్తి}{!!a{derivation}} ఉంటుంది; కనుక ఏదో సంఖ్య $m$కు
$\Prf[\Th{T}](m, \Gn{!G_\Th{T}})$ నిజం. అప్పుడు $\Th{Q}$ వాక్యం
$\OPrf[\Th{T}](\num m, \gn{!G_\Th{T}})$ను \tetoken{నిగమిస్తుంది}{!!{derive}s}. కాబట్టి
$\Th{Q}$ అనేది $\lexists[x][\OPrf[\Th{T}](x,\gn{!G_\Th{T}})]$ను
\tetoken{నిగమిస్తుంది}{!!{derive}s}; నిర్వచనం ప్రకారం అది
$\OProv[\Th{T}](\gn{!G_\Th{T}})$. \olref{eqn:qpf} ప్రకారం $\Th{Q}$
$\lnot !G_\Th{T}$ను \tetoken{నిగమిస్తుంది}{!!{derive}s}; $\Th{T}$ అనేది $\Th{Q}$ను
విస్తరిస్తుంది కాబట్టి $\Th{T}$ కూడా దాన్ని నిగమిస్తుంది. $\Th{T}$ అనేది
$!G_\Th{T}$ను \tetoken{నిగమిస్తే}{!!{derive}s} $\lnot !G_\Th{T}$ను కూడా
\tetoken{నిగమిస్తుందని}{!!{derive}s} చూపాం;
అందువల్ల అది వైరుధ్యమవుతుంది.
\end{proof}

\begin{defn}
\ollabel{thm:oconsis-q}
కింది షరతు నిజమైతే సిద్ధాంతం $\Th{T}$ను \emph{$\omega$-అవైరుధ్యమైనది} అంటాం:
$\lexists[x][!A(x)]$ ఏ వాక్యమైనా, $\Th{T}$ అనేది $\lnot !A(\num 0)$,
$\lnot !A(\num 1)$, $\lnot !A(\num 2)$, \dotsలను \tetoken{నిగమిస్తే}{!!{derive}s}, అప్పుడు
$\lexists[x][!A(x)]$ను $\Th{T}$ నిరూపించదు.
\end{defn}

ప్రతి $\omega$-అవైరుధ్య సిద్ధాంతమూ అవైరుధ్యమైనదని గమనించండి. $\Th{T}$
వైరుధ్యమైతే ప్రతి~$!A$కు $\Th{T} \Proves !A$ అవుతుందనే వాస్తవం నుంచే ఇది
వస్తుంది. ప్రత్యేకించి, $\Th{T}$ వైరుధ్యమైతే, ప్రతి~$n$కు
$\lnot !A(\num n)$ను \tetoken{నిగమిస్తుంది}{!!{derive}s}; అలాగే
$\lexists[x][!A(x)]$ను కూడా \tetoken{నిగమిస్తుంది}{!!{derive}s}. కాబట్టి $\Th{T}$ వైరుధ్యమైతే అది
$\omega$-వైరుధ్యమైనది. విపర్యయ ప్రతిపాదన ప్రకారం, $\Th{T}$
$\omega$-అవైరుధ్యమైనదైతే అది అవైరుధ్యమైనదై ఉండాలి.

\begin{lem}\ollabel{lem:omega-cons-G-unref}
$\Th{T}$ అనేది $\Th{Q}$ను విస్తరించే $\omega$-అవైరుధ్యమైన,
\tetoken{స్వీకృతీకరించదగిన}{!!{axiomatizable}} సిద్ధాంతమైతే,
$\Th{T} \Proves/ \lnot !G_\Th{T}$.
\end{lem}

\begin{proof}
$\Th{T}$ అనేది $\lnot !G_\Th{T}$ను \tetoken{నిగమిస్తే}{!!{derive}s}, అది
$\omega$-వైరుధ్యమైనదని చూపుతాం. $\Th{T}$ అనేది $\lnot !G_\Th{T}$ను
\tetoken{నిగమిస్తుందని}{!!{derive}s} అనుకుందాం. $\Th{T}$ వైరుధ్యమైతే అది
$\omega$-వైరుధ్యమైనది; అప్పుడు పని పూర్తయింది. లేదంటే $\Th{T}$ అవైరుధ్యమైనది;
కాబట్టి \olref{lem:cons-G-unprov} ప్రకారం అది $!G_\Th{T}$ను
\tetoken{నిగమించదు}{!!{derive}}. $\Th{T}$లో $!G_\Th{T}$కు \tetoken{వ్యుత్పత్తి}{!!{derivation}} లేదు కాబట్టి,
$\Th{Q}$ కింది వాటిని \tetoken{నిగమిస్తుంది}{!!{derive}s}:
\[
\lnot \OPrf[\Th{T}](\num 0, \gn{!G_\Th{T}}), \lnot \OPrf[\Th{T}](\num 1,
\gn{!G_\Th{T}}), \lnot \OPrf[\Th{T}](\num 2, \gn{!G_\Th{T}}), \dots
\]
అందువల్ల $\Th{T}$ కూడా వాటిని నిగమిస్తుంది. మరోవైపు, \olref{eqn:qpf}
ప్రకారం $\lnot !G_\Th{T}$ అనేది
$\lexists[x][\OPrf[\Th{T}](x,\gn{!G_\Th{T}})]$కు సమానార్థకం. కాబట్టి
$\Th{T}$ అనేది $\omega$-వైరుధ్యమైనది.
\end{proof}

\begin{prob}
ప్రతి $\omega$-అవైరుధ్య సిద్ధాంతమూ అవైరుధ్యమైనది. దీని విపర్యయం నిజం కాదని,
అంటే అవైరుధ్యమైనప్పటికీ $\omega$-వైరుధ్యమైన సిద్ధాంతాలు ఉన్నాయని చూపండి.
$\Th{Q} \cup \{\lnot !G_\Th{Q}\}$ అవైరుధ్యమైనప్పటికీ
$\omega$-వైరుధ్యమైనదని చూపడం ద్వారా దీన్ని నిరూపించండి.
\end{prob}

\begin{thm}
\ollabel{thm:first-incompleteness} $\Th{T}$ అనేది $\Th{Q}$ను విస్తరించే ఏ
$\omega$-అవైరుధ్యమైన, \tetoken{స్వీకృతీకరించదగిన}{!!{axiomatizable}} సిద్ధాంతమైనా సరే.
అప్పుడు $\Th{T}$ సంపూర్ణం కాదు.
\end{thm}

\begin{proof}
$\Th{T}$ అనేది $\omega$-అవైరుధ్యమైనదైతే, అది అవైరుధ్యమైనది; కాబట్టి
\olref{lem:cons-G-unprov} ప్రకారం $\Th{T} \Proves/ !G_\Th{T}$.
\olref{lem:omega-cons-G-unref} ప్రకారం $\Th{T} \Proves/ \lnot !G_\Th{T}$.
అంటే $\Th{T}$ అనేది $!G_\Th{T}$ను గాని $\lnot !G_\Th{T}$ను గాని
\tetoken{నిగమించదు}{!!{derive}s}; కనుక అది అసంపూర్ణమైనది.
\end{proof}

\end{document}
